% Week X / Chapter template
% Copy into lecture_notes/chapters/weekXX.tex and adapt.

\chapter{Lecture title}
\label{chap:weekXX}

\section{Learning goals}
By the end of this chapter, students should be able to:
\begin{itemize}
  \item 
  \item 
  \item 
\end{itemize}

\section{Why this matters}
Introduce the motivating problem or phenomenon before formalism where possible.

\section{Core ideas}

\subsection{Concept 1}
Explain the concept geometrically and intuitively before introducing notation.

\subsection{Concept 2}
Connect to earlier weeks and make explicit which previous mental model is being extended.

\section{Mathematical formulation}
Introduce only the notation needed to support the conceptual model.

\section{Worked example}
Use a small, inspectable example. Prefer examples that can later be reused in an exercise.

\section{Connection to this week's exercises}

\subsection{Before the experiment}
State what students need to know without revealing the experiment's hidden conclusion.

\subsection{After the experiment}
Provide the formal vocabulary and explanation that consolidates what they observed.

\section{Common misconceptions}
\begin{itemize}
  \item 
  \item 
  \item 
\end{itemize}

\section{Connections forward}
Explain how this week's ideas become prerequisites or intuitions for later weeks.

\section{Summary}
End with a small number of conceptual claims rather than a catalogue of terminology.
